%Master File:lectures.tex

\lesson{Cell Differentiation}
\vspace{-2cm}
\begin{center}
  \includegraphics[height=10cm]{cellDifferentiation.jpg}
\end{center}

\keywords{Random Boolean networks, stem cells, epigenetics}
%%%%%%%%%%%%%%%%%%%%%%% Next Slide %%%%%%%%%%%%%%%%%%%%%%%
\renewcommand{\Outline}{%
\begin{slide}
\section[1]{Outline}

\begin{minipage}{12cm}
  \begin{enumerate}\squeeze
    \outlineitem{Cell Differentiation}{bd}
    \outlineitem{The Life of an Organism}{lifeorg}
    \outlineitem{Epigenetic}{epigenetics}
   \end{enumerate}
\end{minipage}\hfill
\begin{minipage}{10cm}
  \vfill
  \includegraphics[width=10cm]{cell-differentiation}
  \vfill
\end{minipage}
\end{slide}
\addtocounter{outlineitem}{1}
}

\setcounter{outlineitem}{1}
\Outline
\toptarget{firstoutline}

%%%%%%%%%%%%%%%%%%%%%%% Next Slide %%%%%%%%%%%%%%%%%%%%%%%

\begin{slide}
\section{Cell Specialisation}

\begin{PauseHighLight}
  \begin{itemize}
  \item Unlike unicellular organisms most multi-cellular organisms
    exhibit different types of cells\pause
  \item By differentiating, cells can be much better adapted to a
    particular task\pause
  \item In humans there are around 200 different types of cells\pause
  \item Each cell contains the same DNA so cell differentiation involves
    reading a different set of instructions\pause
  \end{itemize}
\end{PauseHighLight}

\end{slide}

%%%%%%%%%%%%%%%%%%%%%%% Next Slide %%%%%%%%%%%%%%%%%%%%%%%

\begin{slide}
\section{Cell Functions}

\begin{PauseHighLight}
  \begin{itemize}
  \item Cell function is determined by the expression of different
    genes\pause
  \item The expression of genes are determined by the transcription
    factors (proteins that interact with the promoter sites either
    exciting or inhibiting the transcription of a gene)\pause
  \item Cells can respond to their environment by the interaction
    between signalling molecules and proteins\pause
  \item Thus, each gene type will respond differently to its environment\pause
  \end{itemize}
\end{PauseHighLight}

\end{slide}

%%%%%%%%%%%%%%%%%%%%%%% Next Slide %%%%%%%%%%%%%%%%%%%%%%%

\begin{slide}
\section[-2]{Regulation by Signal}

\pause
\pb
\pauselevel{=1}
\begin{center}
  \multipdf[width=\linewidth]{geneRegulation}\pause\\
  \multipdf[width=\linewidth]{excitInhib}\pause
\end{center}
\end{slide}

%%%%%%%%%%%%%%%%%%%%%%% Next Slide %%%%%%%%%%%%%%%%%%%%%%%

\begin{slide}
\section{Gene Regulation Networks}

\begin{PauseHighLight}
  \begin{itemize}
  \item Proteins control each other's expression forming complex gene
    regulation networks\pause
  \item The pattern of interaction determines the type of cell\pause
  \item Once the type is determine cells overwhelming stay in that
    type\pause
  \end{itemize}
\end{PauseHighLight}

\end{slide}

%%%%%%%%%%%%%%%%%%%%%%% Next Slide %%%%%%%%%%%%%%%%%%%%%%%

\begin{slide}
\section[-2]{Large Gene Regulation Network}

\begin{center}
  \includegraphics[width=0.75\linewidth]{largeGRN}
\end{center}
\end{slide}

%%%%%%%%%%%%%%%%%%%%%%% Next Slide %%%%%%%%%%%%%%%%%%%%%%%

\begin{slide}
\section{Random Boolean Networks}

\begin{minipage}{0.38\linewidth}
  \includegraphics[width=\linewidth]{randomBooleanFunction}
\end{minipage}\hfil
\begin{minipage}{0.58\linewidth}
  \begin{PauseHighLight}
    \begin{itemize}
    \item To understand gene dynamics researchers (lead by Stuart
      Kauffman) considered the dynamics of random Boolean networks\pause
    \item Each node is connected to $k_{in}$ other nodes uniformly at random\pause
    \item The nodes perform a randomly selected Boolean functions\pause
    \end{itemize}
  \end{PauseHighLight}    
\end{minipage}

\end{slide}

%%%%%%%%%%%%%%%%%%%%%%% Next Slide %%%%%%%%%%%%%%%%%%%%%%%

\begin{slide}
\section[-3]{Dynamics of RBN}

\pause\pb\pauselevel{=1}
\begin{center}
  \multipdf[width=0.7\linewidth]{randomBooleanNetwork}\pause
\end{center}
\end{slide}

%%%%%%%%%%%%%%%%%%%%%%% Next Slide %%%%%%%%%%%%%%%%%%%%%%%

\begin{slide}
\section[-2]{Attractors}

\begin{center}
  \includegraphics[width=0.9\linewidth]{basinsOfAttraction}
\end{center}
\end{slide}

%%%%%%%%%%%%%%%%%%%%%%% Next Slide %%%%%%%%%%%%%%%%%%%%%%%

\begin{slide}
\section{Phase Transition}

\begin{PauseHighLight}
  \begin{itemize}
  \item The behaviour of random Boolean networks shows two phases
    depending on the number of connections and the randomness of the
    functions
    \begin{itemize}
    \item For sparsely connected networks (in-degree $\leq 2$) there are
      a small number of attractors which are reached after a small
      number of iterations\pause
    \item Otherwise there are an exponential number of attractors which
      take exponential time to reach\pause---the ``chaotic'' phase\pause
    \end{itemize}
  \item It has been argued that biology lives on the ``edge of chaos''
    where there are only a few attractors, but the network is
    adaptable\pause
  \end{itemize}
\end{PauseHighLight}

\end{slide}

%%%%%%%%%%%%%%%%%%%%%%% Next Slide %%%%%%%%%%%%%%%%%%%%%%%

\begin{slide}
\section[-3]{Dynamics of RBN in Chaotic Phase}

\pause\pb\pauselevel{=1}
\begin{center}
  \multipdf[width=0.7\linewidth]{chaoticRBN}\pause
\end{center}
\end{slide}

%%%%%%%%%%%%%%%%%%%%%%% Next Slide %%%%%%%%%%%%%%%%%%%%%%%

\begin{slide}
\section{Real Biology}

\begin{PauseHighLight}
  \begin{itemize}
  \item Of course, biological systems are more complex\pause
  \item It is very difficult to turn one type of cell into another\pause
  \item When cells divide they faithfully remember what type they
    are\pause
  \item There are influences beyond the current protein interactions
    that affect a cell's type (epigenetics)\pause
  \end{itemize}
\end{PauseHighLight}

\end{slide}



%%%%%%%%%%%%%%%%%%%%%%% Next Slide %%%%%%%%%%%%%%%%%%%%%%%
\Outline % Life of an Organism
%%%%%%%%%%%%%%%%%%%%%%% Next Slide %%%%%%%%%%%%%%%%%%%%%%%

\begin{slide}
\section{In the Beginning}

\begin{PauseHighLight}
  \begin{itemize}
  \item The life cycle of multi-cellular organisms are complex\pause 
  \item They begin life as a single embryo cell\pause
  \item This cell will divide to form all the specialised cells of the
    organism (and the placenta)\pause
  \item The embryo cell can make any other cell and is known as totipotent\pause
  \item The embryo grows through cell division and quickly
    differentiates into cells of the organism and cells of the
    placenta\pause
  \end{itemize}
\end{PauseHighLight}

\end{slide}

%%%%%%%%%%%%%%%%%%%%%%% Next Slide %%%%%%%%%%%%%%%%%%%%%%%

\begin{slide}
\section{Embryonic Stem Cells}

\begin{PauseHighLight}
  \begin{itemize}
  \item The cells in the organism are no longer able to form a placenta,
    but can produce any other cell in the body\pause
  \item These \emph{embryonic stem cells} are said to be
    \emph{pluripotent}\pause
  \item The embryonic stem cells can change to more specialised stem
    cells which are said to be \textit{multipotent}\pause
  \item These will eventually specialise to the normal cells of the body
    known as the \emph{soma}\pause
  \end{itemize}
\end{PauseHighLight}

\end{slide}

%%%%%%%%%%%%%%%%%%%%%%% Next Slide %%%%%%%%%%%%%%%%%%%%%%%

\begin{slide}
\section{Specialisation}

\begin{center}
  \includegraphics[width=\linewidth]{specialisation}
\end{center}
\end{slide}

%%%%%%%%%%%%%%%%%%%%%%% Next Slide %%%%%%%%%%%%%%%%%%%%%%%

\begin{slide}
\section{Causes of Specialisation}

\begin{minipage}{0.28\linewidth}
  \begin{center}
  \includegraphics[width=0.5\linewidth]{asymmetricCellDivision}\\
  \vspace*{1cm}
  \includegraphics[width=\linewidth]{signallingAtReceptors}    
  \end{center}
\end{minipage}\hfil
\begin{minipage}{0.68\linewidth}
\begin{PauseHighLight}
  \begin{itemize}
  \item There are a number of routes to specialisation\pause
  \item During the development of the embryo some transcription factors
    aggregate, at cell division this leads to differentiation in the
    cell function\pause
  \item Other cells are known to specialise depending on signalling
    molecules that occur at receptors\pause
  \end{itemize}
\end{PauseHighLight}
  
\end{minipage}


\end{slide}


%%%%%%%%%%%%%%%%%%%%%%% Next Slide %%%%%%%%%%%%%%%%%%%%%%%
\Outline % Epigenetics
%%%%%%%%%%%%%%%%%%%%%%% Next Slide %%%%%%%%%%%%%%%%%%%%%%%

\begin{slide}
\section{Cell Division}

\begin{PauseHighLight}
  \begin{itemize}
  \item Cells in the body constantly divide\pause
  \item They all carry the same DNA code\pause
  \item But, the DNA is often modified by methylation and
    acetylation\pause
  \item These can be facilitated, e.g. by proteins called \textit{DNA
      methyltransferase} (DNMT)\pause
  \end{itemize}
\end{PauseHighLight}

\end{slide}

%%%%%%%%%%%%%%%%%%%%%%% Next Slide %%%%%%%%%%%%%%%%%%%%%%%

\begin{slide}
\section{DNA Methylation}

\begin{center}
  \includegraphics[width=\linewidth]{methylationPic}
\end{center}
\end{slide}

%%%%%%%%%%%%%%%%%%%%%%% Next Slide %%%%%%%%%%%%%%%%%%%%%%%

\begin{slide}
\section[-2]{Histone Acetylation}

\begin{center}
  \includegraphics[width=0.95\linewidth]{acetylation}
\end{center}
\end{slide}

%%%%%%%%%%%%%%%%%%%%%%% Next Slide %%%%%%%%%%%%%%%%%%%%%%%

\begin{slide}
\section{Epigenetics}

\begin{PauseHighLight}
  \begin{itemize}
  \item Acetylation suppresses the expression of genes\pause
  \item DNA methylation can suppress or activate the expression of
    genes\pause
  \item Stem cells have more methylated sites than somatic (normal)
    cells, showing that this is an important mechanism for ensuring cell
    specialisation is maintained during cell division\pause
  \item The methylation and acetylation is modified by proteins\pause
  \item Epigenetics can be seen as another component in a complex
    Boolean network\pause
  \end{itemize}
\end{PauseHighLight}

\end{slide}

%%%%%%%%%%%%%%%%%%%%%%% Next Slide %%%%%%%%%%%%%%%%%%%%%%%

\begin{slide}
\section{Conclusion}

\begin{PauseHighLight}
  \begin{itemize}
  \item Cell differentiation is a complex biological process that is
    essential to multi-cellular organism\pause
  \item Simplified models of Boolean networks suggests that a small
    number of attractors might naturally occur in an interacting
    network\pause
  \item Whether this ``explains'' multi-cellular organisms or their
    origin is a matter of debate\pause
  \item Specialisation happens through a hierarchical series of
    steps\pause
  \item Epigenetics seem to be an important component in determining
    which proteins are expressed\pause
  \end{itemize}
\end{PauseHighLight}

\end{slide}

